\section{Discussion}
\label{sec:discussion}

\subsection{Availability did not provide sufficient support for this benchmark}
The study separates access to a product from support for a particular evaluation. Official files were available, 128 sources were processed, and their extracts and accounting records could be checked. Nevertheless, the selected stations and period provided only 150 of the 400 required opportunities, with most station--block cells below their local threshold. The original performance experiment could not proceed under its declared rules.

This result adds a design-specific question to the concerns addressed by analysis-ready data and quality-assurance frameworks: after the population, matching criteria, and reporting cells have been fixed, how much usable support remains for the intended comparison? Consistent processing and traceable quality information help make data usable, but do not supply the answer to that question~\cite{dwyer2018,qa4eo2010}. Nor do published coastal ASCAT validation results determine the opportunity counts for this selection~\cite{verhoef2012,lindsley2016,soisuvarn2023}. Nothing in the present counts establishes that the product is inaccurate or unsuitable for other applications.

\subsection{The stopping decision and its replay answer different questions}
The capacity gate prevented the workflow from continuing to outcomes when its support requirements were unmet. Replacing stations, extending the interval, relaxing a criterion, or choosing individual alternate files could define a useful new experiment. Those changes would not show that the original design had passed. Keeping the complete alternate rounds and stopping after their exhaustion made the distinction explicit.

The replay addressed whether that decision and its supporting arithmetic could be reconstructed. Its successful checks improve the inspectability of the recorded failure: duplicated opportunities, inflation by neighborhood pixels, missing final records, or an unexecuted prescribed alternate round do not explain the reproduced count. They do not establish the physical cause of the shortfall, independently validate the extraction software, or confer external reproducibility. Separating the capacity decision from record integrity avoids treating successful data engineering as evidence that a scientific sampling requirement has been met.

\subsection{Thresholds should follow the intended analysis}
The 400-total and eight-per-cell requirements were inherited from the original workflow. They allow losses of 20\% overall and 25\% within a cell before the downstream minima of 320 and six would be crossed. These calculations explain the relationship between the two stages; they do not justify either set of minima through statistical power or target precision.

A successor study should specify its estimand, independent analysis unit, target precision or power, relevant heterogeneity, and expected retention between predictor screening and outcome analysis. The station and file populations, local support requirements, and acquisition budget can then be set together. This would make a failed capacity check informative about a justified scientific design, rather than only about compliance with an inherited configuration. The current negative result remains useful because it identifies the support that was actually available before any reported outcome access.

\subsection{What can and cannot be learned from sensitivity calculations}
The terminal counts provide one simple upper bound. Only 35 evaluation rows failed the two-candidate rule. Even if all 35 became valid after relaxing that rule, and no other requirement changed, the total could be no greater than 185, still below 400. This is an arithmetic ceiling, not a completed one-candidate experiment. Determining the actual gain would require new selection semantics and reprocessing.

The retained material is insufficient for an analogous 15- or 20-km spatial calculation. Failed rows do not preserve the nearest center outside the 12.5-km tolerance, and the 125 downloaded raw files were deleted after extraction. A future outcome-blind study could preregister a distance-by-candidate-rule grid, reacquire the source files, and retain nearest-center distances for every terminal state. It should report the resulting counts under the new design rather than merge them into the present result.

The acquisition trace also leaves a separate optimization question open. Complete rounds maintained representation across the eight strata and prevented selection of individual files after their opportunity yields were seen. Comparing that policy with progressive activation or forecasts of orbital and buoy sampling would require a larger prespecified candidate universe, explicit objectives, and the same budget and stopping conditions. The five realized points here cannot establish which policy would be more efficient.

\subsection{Making the evidence usable beyond its original repository}
A reusable readiness record needs both the decision rule and enough information to reconstruct its inputs. The population, interval, counting unit, deduplication procedure, thresholds, alternate order, budget, outcome-access boundary, terminal definitions, digests, and stop state should travel with the result. The source manuscript reports a project-level JSON Schema for these fields; it is an implementation aid, not a community standard.

Existing infrastructures can represent complementary parts of such a release. STAC can describe the candidate assets and their spatial and temporal extents~\cite{stac2025}. RO-Crate and Workflow Run RO-Crate can link plans, software, inputs, outputs, and execution records~\cite{rocrate2022,leo2024}, while BagIt provides a payload-manifest approach to fixity~\cite{bagit2018}. Pangeo-style repositories address access to scientific arrays and scalable analysis~\cite{abernathey2021}. These are possible mappings, not standards conformance demonstrated by this implementation.

The remaining release work is concrete: resolve the stale locator, retain probe-level receipts in future acquisitions, document redistribution rights, provide a locked runtime and checksums, and deposit the permitted artifacts in a durable versioned archive. An independently implemented replay would then test reproducibility beyond the original software context, consistent with the broader reproducibility goals of remote sensing research~\cite{kedron2022}. Until those steps are completed, the evidence supports the reported design decision and its internal reconstruction.
